\begin{center}
    {\LARGE \textbf{2013无机化学}} \\[2ex]
    {\large 时量180分钟 \quad 满分100分}
\end{center}

\vspace{0.5cm}
\noindent\textbf{注意事项:}
\begin{enumerate}[label=\arabic*., parsep=0pt, itemsep=0pt]
    \item 答题前,考生先将自己的姓名、学号、考场号、座位号填写在答题卡上,将条形码准确粘贴在考生信息条形码粘贴区。
    \item 考试结束后,将本试卷与答题纸一并收回。
    \item 回答各个问题时,将答案写在答题纸上,写在本试卷上无效。
\end{enumerate}

\vspace{0.5cm}
\noindent\textbf{一、完成并配平下列反应方程式(15 pts.)}

\begin{enumerate}[label=\arabic*., start=1, itemsep=1.5ex]
    \item 无氧条件下用\ce{Zn}粉还原亚硫酸氢钠溶液。
    \item 将\ce{CO}通入\ce{CuCl}的酸性溶液中。
    \item 铜溶于氰化钠溶液。
    \item 向盛有\ce{HgCl2}溶液的试管中滴加\ce{SnCl2}溶液。
    \item 单质硅溶于氢氧化钠的浓溶液。
\end{enumerate}

\vspace{0.5cm}
\noindent\textbf{二、回答下列各题(83 pts.)}

\begin{enumerate}[label=\arabic*., start=1, itemsep=1.5ex]
    \item 卤素分子\ce{F2}、\ce{Cl2}、\ce{Br2}、\ce{I2}的解离能分别为155、240、190、149\ \ce{kJ\cdot mol^{-1}}，简要说明为什么\ce{F2}的解离能小于\ce{Cl2}、\ce{Br2}、而和\ce{I2}相近。
    \item 画出\ce{OF}分子轨道图，分别写出\ce{OF}、\ce{OF+}、\ce{OF-}的分子轨道排布式，并比较各物质的稳定性。
    \item 解释下列实验现象，并简要说明原因。
    \begin{enumerate}[label=(\arabic*), itemsep=1ex]
        \item 稀释\ce{CuCl2}浓溶液时，溶液的颜色依次是:黄色，黄绿，绿色，蓝绿，蓝色。
        \item 硼砂与\ce{CoCl2}共熔得到蓝色产物，硼砂与\ce{CrCl3}共熔得到绿色产物。
        \item 向\ce{CoCl2}溶液中加稀\ce{KSCN}溶液，溶液不变蓝;再加入乙醚则有机层变蓝。
        \item 向\ce{K2MnO4}溶液中缓慢通入\ce{NO2}，先有棕黑色沉淀生成，\ce{NO2}过量则沉淀消失得到无色溶液。
    \end{enumerate}
    \item 写出原子序数为31、44和54的元素符号、名称和价层电子的构型。
    \item 简述\ce{OSF2}、\ce{OSCl2}和\ce{OSBr2}分子中\ce{S-O}键强度的变化规律，并解释原因。
    \item \ce{Co^{3+}}和\ce{Fe^{3+}}的电荷相同，半径也相近(分别为63\ \ce{pm}和64\ \ce{pm})，为什么\ce{Co^{3+}}配位化合物比\ce{Fe^{3+}}配位化合物稳定?
    \item 用两种方法鉴别下列各对物质。
    \begin{enumerate}[label=(\arabic*), itemsep=1ex]
        \item \ce{MnO2}和\ce{CuO}
        \item \ce{Na3PO4}和\ce{NaAsO4}
    \end{enumerate}
    \item 指出下列配离子中心的杂化类型、几何构型和磁矩。
    \begin{enumerate}[label=(\arabic*), itemsep=1ex]
        \item \ce{[Pt(CN)4]^{2-}}
        \item \ce{[Mn(CN)6]^{4-}}
        \item \ce{Fe(CO)5}
        \item \ce{[Co(SCN)4]^{2-}}
    \end{enumerate}
    \item 在酸性溶液中，\ce{KBrO3}能把\ce{KI}氧化成\ce{I2}和\ce{KIO3}，本身可被还原为\ce{Br2}、\ce{Br-}；而\ce{KIO3}和\ce{KBr}反应生成\ce{I2}和\ce{Br2}，\ce{KIO3}和\ce{KI}反应生成\ce{I2}，现于酸性溶液中混合等物质的量的\ce{KBrO3}和\ce{KI}，生成哪些氧化还原产物。它们的物质的量的比是多少?
    \item 已知\ce{Tl}的元素电势图
    $E^\ominus/\text{V} \quad 
    \ce{Tl^{3+}} 
    \xline{\raisebox{0.5ex}{$+1.25$}} 
    \ce{Tl^{+}} 
    \xline{\raisebox{0.5ex}{$-0.3363$}} 
    \ce{Tl}$

    氧化还原反应: \ce{2Tl + Tl^{3+} = 3Tl^{+}}，能够设计成几个原电池?写出电极反应和电池符号，计算各原电池的电动势$E^\ominus$。
\end{enumerate}

\vspace{0.5cm}
\noindent\textbf{三、制备与分离（28 pts.）}

\begin{enumerate}[label=\arabic*., start=1, itemsep=1.5ex]
    \item 说明工业上以\ce{Ca3(PO4)2}为原料制取单质磷的方法。
    \item 以铬绿(\ce{Cr2O3})为原料制备重铬酸钾，写出相关的反应方程式。
    \item 如何除去\ce{KCl}溶液中含有的少量的\ce{K2SO4}杂质?写出相关反应方程式。
    \item 设计方案将混合溶液中的\ce{Al^{3+}}、\ce{Zn^{2+}}、\ce{Fe^{3+}}、\ce{Ni^{2+}}、\ce{Cu^{2+}}分离。
\end{enumerate}

\vspace{0.5cm}
\noindent\textbf{四、推断（24 pts.）}

\begin{enumerate}[label=\arabic*., start=1, itemsep=1.5ex]
    \item 白色晶体\ce{A}加热后有红棕色气体\ce{B}生成，向\ce{A}的水溶液中加入盐酸有白色沉淀\ce{C}生成，盐酸过量时沉淀\ce{C}溶解得到\ce{D}。若向\ce{A}的水溶液中加入\ce{KI}溶液，则有黄色沉淀\ce{E}生成。向\ce{A}的水溶液中滴加\ce{NaOH}溶液有白色沉淀\ce{F}生成，\ce{NaOH}过量则沉淀\ce{F}溶解。将气体\ce{B}通入水中生成\ce{G}，并放出无色气体\ce{H}。沉淀\ce{F}溶于适量\ce{G}中形成的溶液结晶后可以得到\ce{A}。
    请给出\ce{A}、\ce{B}、\ce{C}、\ce{D}、\ce{E}、\ce{F}、\ce{G}和\ce{H}的化学式，并用化学反应方程式表示各过程。
    \item 绿色水合盐\ce{A}在一定温度下脱水得到白色固体\ce{B}。\ce{B}在更高的温度下分解生成红棕色固体\ce{C}和气体\ce{D}，气体\ce{D}能使\ce{KMnO4}溶液褪色。\ce{C}溶于盐酸后加入\ce{KSCN}溶液得红色溶液\ce{E}。向溶液\ce{E}中加入\ce{NaOH}溶液生成棕色沉淀\ce{F}。\ce{F}溶于\ce{KHC2O4}溶液得到黄色溶液，经蒸发、浓缩后冷却，析出绿色水合晶体\ce{G}。\ce{F}溶于\ce{KHSO3}溶液后，将溶液浓缩则有\ce{A}析出。
    请给出\ce{A}、\ce{B}、\ce{C}、\ce{D}、\ce{E}、\ce{F}、\ce{G}的化学式及相关反应的化学方程式。
    \item 化合物\ce{A}为白色固体，不溶于水。\ce{A}受热剧烈分解，生成固体\ce{B}和气体\ce{C}。固体\ce{B}不溶于水或盐酸，但溶于热的稀硝酸得到无色溶液\ce{D}和无色气体\ce{E}。\ce{E}在空气中变棕色。向溶液\ce{D}中加入盐酸时得到白色沉淀\ce{F}。气体\ce{C}与普通试剂不起反应，但与热的金属镁作用生成白色固体\ce{G}。\ce{G}与水作用得白色沉淀\ce{H}及气体\ce{I}。\ce{I}能使湿润的红色石蕊试纸变蓝。\ce{H}可溶于稀硫酸得到溶液\ce{J}。化合物\ce{A}以\ce{H2S}溶液处理时得黑色沉淀\ce{K}、无色溶液\ce{L}和气体\ce{C}。过滤后，固体\ce{K}溶于浓硝酸得气体\ce{E}、黄色沉淀\ce{M}和溶液\ce{D}。用\ce{NaOH}溶液处理滤液\ce{L}又得气体\ce{I}。
    问\ce{A}、\ce{B}、\ce{C}、\ce{D}、\ce{E}、\ce{F}、\ce{G}、\ce{H}、\ce{I}、\ce{J}、\ce{K}、\ce{L}、\ce{M}各为何物，写出由\ce{B}生成\ce{D}和\ce{E}，及由\ce{A}生成\ce{K}、\ce{L}和\ce{C}的反应方程式？
\end{enumerate}